% 论文层级类图内嵌效果演示（hierarchy.md 指导）：组织架构（含虚线汇报）+ 思维导图 + 概念图
% 编译：xelatex -interaction=nonstopmode paper_embedded_hierarchy.tex（两遍，交叉引用）
% fontset=none：禁止 ctex 自行选择 CJK 字体（它在 Linux 上会去找 Fandol 并可能失败），
% 中英文字体统一由下面的 \input{.../fonts-serif.tex} 决定（见 references/common.md「字体与文字」§1）
\documentclass[12pt,fontset=none]{ctexart}
\ctexset{fontset=none}
\usepackage[a4paper, margin=2.6cm]{geometry}
% ============================================================
% pretty-charts TikZ/pgfplots 导言区模板 —— academic 档（T1 出版级）
% 编译方式：XeLaTeX。business/showcase 档改字号与 cycle list（见文件末尾注释）
% ============================================================
\usepackage{pgfplots}
\pgfplotsset{compat=1.18}
\usetikzlibrary{arrows.meta, positioning, calc}

% ---- 图表字体：中文黑体 + 西文 Arial（XeLaTeX；正文仍可用全局宋体+TNR，互不影响）----
% 注意：Windows 的 Noto Sans SC 多为可变字体(VF)，xdvipdfmx 无法嵌入（fatal: Invalid font），
% 故 TikZ 导出 PDF 用 Microsoft YaHei；装有静态版思源黑体(Source Han Sans SC)的机器可换回。
\usepackage{fontspec}
\usepackage{xeCJK}
\IfFontExistsTF{Microsoft YaHei}
  {\newCJKfontfamily\chartcjk{Microsoft YaHei}}% 静态字体，可嵌入 PDF
  {\IfFontExistsTF{Noto Sans CJK SC}
     {\newCJKfontfamily\chartcjk{Noto Sans CJK SC}}% Linux（fonts-noto-cjk）
     {\IfFontExistsTF{FandolHei}
        {\newCJKfontfamily\chartcjk{FandolHei}}% 兜底：TeX Live 自带
        {\PackageWarning{pretty-charts}{未找到中文黑体，图表文字可能显示异常}}}}
\IfFontExistsTF{Arial}
  {\newcommand{\chartfont}[1]{{\chartcjk\fontspec{Arial}#1}}}
  {\newcommand{\chartfont}[1]{{\chartcjk\fontspec{TeX Gyre Heros}#1}}}

% ---- 色板：academic（Okabe-Ito 重排，与 references/assets/palettes/academic.json 一致）----
\definecolor{pcBlue}{HTML}{0072B2}
\definecolor{pcOrange}{HTML}{D55E00}
\definecolor{pcGreen}{HTML}{009E73}
\definecolor{pcMagenta}{HTML}{CC79A7}
\definecolor{pcYellow}{HTML}{E69F00}
\definecolor{pcCyan}{HTML}{56B4E9}
\definecolor{pcGray}{HTML}{999999}
\definecolor{pcGrid}{HTML}{CCCCCC}
\definecolor{pcAxis}{HTML}{333333}
\definecolor{pcTick}{HTML}{4D4D4D}
\definecolor{pcTitleBlack}{HTML}{1A1A1A} % = palettes/academic.json text.title
\definecolor{pcNeutral}{HTML}{B3B3B3}    % = palettes/academic.json neutral（中性灰，spec/color.md §4.4）

% ---- 全局轴风格：去顶右刺、浅色 y 网格、字号层级 ----
\pgfplotsset{
  every axis/.append style={
    axis line style={pcAxis, line width=0.4pt},
    label style={font=\small, text=pcAxis},
    tick label style={font=\footnotesize, text=pcTick},
    title style={font=\small\bfseries, text=pcTitleBlack},
    legend style={font=\footnotesize, draw=none, fill=none, nodes={scale=0.9, transform shape}},
    grid=major,
    major grid style={pcGrid, line width=0.4pt},
    axis x line*=bottom,
    axis y line*=left,
    tick style={pcTick, line width=0.4pt},
  },
  every axis plot/.append style={line width=0.9pt},
  /pgfplots/error bars/error bar style={line width=0.5pt},
}

% ---- 色板循环列表 ----
\pgfplotscreateplotcyclelist{pc academic}{
  pcBlue\\ pcOrange\\ pcGreen\\ pcMagenta\\ pcYellow\\ pcCyan\\ pcGray\\
}

% ---- 用法示例 ----
% \begin{tikzpicture}
%   \begin{axis}[cycle list name=pc academic, xlabel={年份}, ylabel={销量}, title={标题}]
%     \addplot coordinates {(1,2) (2,4) (3,3)};
%   \end{axis}
% \end{tikzpicture}

% ============================================================
% business 档（T2）差异：字号放大一档（label \normalsize / tick \small），
% cycle list 换 4E79A7/F28E2B/E15759/76B7B2/59A14F/EDC948/B07AA1/FF9DA7/9C755F/BAB0AC
% showcase 档（T3）：字号再放大（label \large / tick \normalsize），line width=1.4pt，
% cycle list 换 0077BB/EE7733/009988/EE3377/33BBEE/CC3311/BBBBBB
% ============================================================

% ============================================================
% pretty-charts TikZ 流程图预置样式
% 论文版与 PPT 版共用同一套节点形状，仅配色/线宽不同：
%   \pcPaperMode  ← 论文（默认）：黑白灰，无填充色
%   \pcPPTMode    ← 演示：showcase 色板彩色，粗线大字
% 用法：先 \input{references/assets/tikz/preamble.tex}，再 \input 本文件，正文选模式后作图。
% 依赖库已在此加载：shapes.geometric, arrows.meta, positioning
% ============================================================
\usetikzlibrary{shapes.geometric, arrows.meta, positioning, fit, backgrounds}

% ---- 论文版（黑白无底色；字体由文档设：中文宋体 + 西文 Times New Roman）----
\newcommand{\pcPaperMode}{%
  \def\pcNodeLine{black!75}\def\pcNodeFill{white}\def\pcNodeText{black}%
  \def\pcKeyFill{white}\def\pcKeyLine{black}\def\pcKeyText{black}%
  \def\pcKeyWidth{1pt}% 关键节点靠粗边框区分，不用灰底
  \def\pcSubLine{black!40}\def\pcSubText{black!65}%
  \def\pcArrow{black!80}\def\pcLabel{black!70}%
  \def\pcNodeFont{\small}\def\pcFlowWidth{0.6pt}\def\pcStealth{2.2mm}}

% ---- PPT 版（showcase 色板：蓝主色、橙强调；字体建议中文黑体 + 西文 Arial）----
\newcommand{\pcPPTMode}{%
  \definecolor{pcPptBlue}{HTML}{0077BB}%
  \definecolor{pcPptOrange}{HTML}{EE7733}%
  \def\pcNodeLine{pcPptBlue!80}\def\pcNodeFill{pcPptBlue!10}\def\pcNodeText{black}%
  \def\pcKeyFill{pcPptBlue}\def\pcKeyLine{pcPptBlue}\def\pcKeyText{white}%
  \def\pcKeyWidth{\pcFlowWidth}%
  \def\pcSubLine{pcPptBlue!40}\def\pcSubText{black!65}%
  \def\pcArrow{black!70}\def\pcLabel{pcPptOrange!90!black}%
  \def\pcNodeFont{\normalsize}\def\pcFlowWidth{1.1pt}\def\pcStealth{3mm}}

\pcPaperMode % 默认论文版；PPT 场景在文档里改调 \pcPPTMode

% ---- 节点形状（形状语义与 references/diagram/flowchart.md 一致）----
\tikzset{
  % 起止节点：圆角胶囊
  pc start/.style={draw=\pcNodeLine, fill=\pcNodeFill, rounded corners=10pt,
                   minimum height=0.72cm, inner sep=7pt,
                   font=\pcNodeFont, text=\pcNodeText, align=center},
  % 处理：直角矩形
  pc process/.style={draw=\pcNodeLine, fill=\pcNodeFill, rounded corners=1.5pt,
                     minimum width=2.2cm, minimum height=0.75cm, inner sep=5pt,
                     font=\pcNodeFont, text=\pcNodeText, align=center},
  % 关键节点（关键路径）：论文版加粗边框（无底色），PPT 版主色实心白字
  pc key/.style={pc process, draw=\pcKeyLine, fill=\pcKeyFill, text=\pcKeyText,
                 line width=\pcKeyWidth, font=\pcNodeFont},
  % 判断：菱形
  pc decision/.style={shape=diamond, aspect=1.7, inner sep=1.5pt,
                      draw=\pcNodeLine, fill=\pcNodeFill,
                      font=\pcNodeFont, text=\pcNodeText, align=center},
  % 输入/输出：平行四边形
  pc io/.style={shape=trapezium, trapezium left angle=75, trapezium right angle=105,
                inner sep=4pt, minimum height=0.72cm,
                draw=\pcNodeLine, fill=\pcNodeFill,
                font=\pcNodeFont, text=\pcNodeText, align=center},
  % 数据库：圆柱
  pc db/.style={shape=cylinder, shape border rotate=90, cylinder uses custom fill,
                aspect=0.35, inner sep=2pt, minimum width=1.9cm, minimum height=1.0cm,
                draw=\pcNodeLine, fill=\pcNodeFill,
                font=\pcNodeFont, text=\pcNodeText, align=center},
  % 阶段/边界容器：虚线框
  pc sub/.style={draw=\pcSubLine, dashed, rounded corners=3pt, inner sep=8pt},
  pc sub label/.style={font=\footnotesize, text=\pcSubText},
}

% ---- 箭头与标注 ----
\tikzset{
  pc flow/.style={-{Stealth[length=\pcStealth]}, line width=\pcFlowWidth, draw=\pcArrow},
  pc flow back/.style={pc flow, dashed},                 % 回流/反馈：虚线
  pc label/.style={font=\footnotesize, text=\pcLabel, fill=white, inner sep=1.5pt},
}

% ---- 用法示例 ----
% \begin{tikzpicture}[node distance=8mm and 10mm]
%   \node[pc start]                (s) {开始};
%   \node[pc process, below=of s]  (a) {参数校验};
%   \node[pc decision, below=of a] (b) {库存充足？};
%   \draw[pc flow] (s) -- (a);
%   \draw[pc flow] (a) -- node[pc label, right=2pt]{通过} (b);
% \end{tikzpicture}

\usepackage{caption}
\captionsetup{font=small, labelsep=quad}
\setcounter{section}{2}
\renewcommand{\thefigure}{\thesection.\arabic{figure}}
% ============================================================
% pretty-charts 字体层：衬线搭配（宋体 + Times New Roman）
% 用法：\input{.../references/assets/tikz/fonts-serif.tex}
%
% 为什么单独成文件：示例此前各自写 \setmainfont{Times New Roman}
% + \setCJKmainfont{SimSun}，这两个字体只存在于 Windows，同一份示例在
% Linux/macOS 上无法编译（CI 也就无从编译）。这里按"字体是否存在"逐级回退：
%   Windows  → Times New Roman + SimSun（与原行为完全一致）
%   Linux    → TeX Gyre Termes + Noto Serif CJK SC（apt install fonts-noto-cjk-extra）
%   兜底     → TeX Gyre Termes + FandolSong（TeX Live 自带）
%   再兜底   → 只发警告、不设字体，保证编译不因缺字体而失败（中文可能显示异常）
% \IfFontExistsTF 由 fontspec 提供（preamble.tex 已加载 fontspec）。
% ============================================================
\IfFontExistsTF{Times New Roman}
  {\setmainfont{Times New Roman}}          % Windows / macOS
  {\IfFontExistsTF{TeX Gyre Termes}
     {\setmainfont{TeX Gyre Termes}}        % Linux 需 apt install fonts-texgyre
     {\PackageWarning{pretty-charts}{未找到西文衬线字体（Times New Roman / TeX Gyre %
       Termes），改用 XeLaTeX 默认西文字体；如需一致字形请安装 fonts-texgyre}}}
\IfFontExistsTF{SimSun}
  {\setCJKmainfont{SimSun}}                % 中文宋体（Windows）
  {\IfFontExistsTF{Noto Serif CJK SC}
     {\setCJKmainfont{Noto Serif CJK SC}}   % Linux（fonts-noto-cjk-extra）
     {\IfFontExistsTF{FandolSong}
        {\setCJKmainfont{FandolSong}}        % 兜底：TeX Live 自带，无需系统字体包
        {\PackageWarning{pretty-charts}{未找到中文衬线字体（已尝试 SimSun / %
          Noto Serif CJK SC / FandolSong）；中文可能显示异常。%
          请安装 fonts-noto-cjk-extra，或改用 fonts-sans.tex}}}}

\usetikzlibrary{mindmap}

\begin{document}
\section{系统设计与实现}

研发团队的组织设计遵循"职能线与质量线双轨并行"的原则：各部门实线汇报至总经理，
质量条线以虚线汇报至董事会审计委员会，保证测试环节的独立性。组织结构、效能提升
思路与核心概念间的关系分别如图~\ref{fig:org}、图~\ref{fig:mind}~与图~\ref{fig:concept}~所示。

% ================= 图 3.1 组织架构图 =================
\begin{figure}[htbp]
  \centering
  \begin{tikzpicture}
    % ---- 节点 ----
    \node[pc process] (board) at (0,1.7)    {董事会};
    \node[pc process] (audit) at (4.6,1.7)  {审计委员会};
    \node[pc key]     (ceo)   at (0,0)      {总经理};
    \node[pc process] (rd)    at (-5.2,-1.9){研发中心};
    \node[pc process] (mkt)   at (0.7,-1.9) {市场部};
    \node[pc process] (hr)    at (4.2,-1.9) {人事部};
    \node[pc process] (fin)   at (6.8,-1.9) {财务部};
    \node[pc process] (fe)    at (-7.2,-3.9){前端组};
    \node[pc process] (be)    at (-5.2,-3.9){后端组};
    \node[pc process] (qa)    at (-3.2,-3.9){测试组};
    \node[pc process] (dom)   at (-0.2,-3.9){国内部};
    \node[pc process] (ovs)   at (1.6,-3.9) {海外部};
    \node[pc process] (rec)   at (3.4,-3.9) {招聘组};
    \node[pc process] (trn)   at (5.0,-3.9) {培训组};

    % ---- 治理层：实线=直接汇报，虚线=虚线汇报 ----
    \draw[pc flow] (board.south) -- (ceo.north);
    \draw[pc flow back] (audit.west) -- node[pc label, above=1pt]{虚线汇报} (board.east);

    % ---- 一级汇报：竖干 + 横轨 + 落入各部门 ----
    \draw[pc flow] (ceo.south) -- (0,-0.95);
    \draw[draw=black!80, line width=0.6pt] (-5.2,-0.95) -- (6.8,-0.95);
    \foreach \x/\d in {-5.2/rd, 0.7/mkt, 4.2/hr, 6.8/fin}
      \draw[pc flow] (\x,-0.95) -- (\d.north);

    % ---- 二级汇报 ----
    \draw (rd.south) -- (-5.2,-2.9);
    \draw[draw=black!80, line width=0.6pt] (-7.2,-2.9) -- (-3.2,-2.9);
    \foreach \x/\d in {-7.2/fe, -5.2/be, -3.2/qa}
      \draw[pc flow] (\x,-2.9) -- (\d.north);
    \draw (mkt.south) -- (0.7,-2.9);
    \draw[draw=black!80, line width=0.6pt] (-0.2,-2.9) -- (1.6,-2.9);
    \foreach \x/\d in {-0.2/dom, 1.6/ovs}
      \draw[pc flow] (\x,-2.9) -- (\d.north);
    \draw (hr.south) -- (4.2,-2.9);
    \draw[draw=black!80, line width=0.6pt] (3.4,-2.9) -- (5.0,-2.9);
    \foreach \x/\d in {3.4/rec, 5.0/trn}
      \draw[pc flow] (\x,-2.9) -- (\d.north);
  \end{tikzpicture}
  \caption{研发团队组织架构。实线为直接汇报关系，虚线为审计条线的虚线汇报；
  加粗节点为关键管理岗位。}
  \label{fig:org}
\end{figure}

图~\ref{fig:org}~中，测试组除实线汇报研发中心外，其质量结论同时虚线汇报审计
委员会——双轨设计使质量评估独立于交付节奏。组织架构图遵循"一人一框、同层
互斥"的规范，跨部门协作不在此图中表达（那是流程图的职责）。

% ================= 图 3.2 思维导图 =================
\begin{figure}[htbp]
  \centering
  \begin{tikzpicture}[mindmap, grow cyclic,
    every node/.style=concept, concept color=black!8,
    root concept/.append style={concept color=black!18, font=\small\bfseries,
                                minimum size=2.8cm, text width=2.4cm},
    level 1/.append style={level distance=4.3cm, sibling angle=90,
                           font=\footnotesize, minimum size=1.8cm, text width=1.6cm},
    level 2/.append style={level distance=3.0cm, sibling angle=60,
                           font=\footnotesize, minimum size=1.35cm, text width=1.15cm}]
    \node[root concept]{研发效能}
      child {node{流程}
        child {node{需求评审}}
        child {node{自动化测试}}
        child {node{持续集成}}}
      child {node{工具}
        child {node{代码审查}}
        child {node{效能平台}}}
      child {node{组织}
        child {node{技能培训}}
        child {node{团队协作}}}
      child {node{度量}
        child {node{交付周期}}
        child {node{缺陷密度}}};
  \end{tikzpicture}
  \caption{研发效能提升思维导图（一级分支 4 个，层级 ≤3）。}
  \label{fig:mind}
\end{figure}

% ================= 图 3.3 概念图 =================
\begin{figure}[htbp]
  \centering
  \begin{tikzpicture}
    \node[pc key]     (dtype) at (-4.4,0.6)  {数据类型};
    \node[pc process] (chan)  at (-0.4,2.4)  {视觉通道};
    \node[pc process] (task)  at (3.8,1.7)   {读者任务};
    \node[pc key]     (sel)   at (1.0,-1.4)  {图型选择};
    \node[pc process] (eff)   at (-4.6,-2.0) {编码效率};
    \node[pc process] (mis)   at (4.2,-2.8)  {防误导规范};

    \draw[pc flow] (dtype.east) -- node[pc label, above=1pt]{决定} (chan.west);
    \draw[pc flow] (chan.east) -- node[pc label, above=1pt]{服务于} (task.west);
    \draw[pc flow] (task.south) -- node[pc label, right=2pt]{驱动} (sel.north east);
    \draw[pc flow] (eff.north) -- node[pc label, above=1pt, pos=0.4]{决定} (chan.south west);
    \draw[pc flow] (sel.east) -- node[pc label, above=1pt]{遵循} (mis.west);
    \draw[pc flow] (dtype.south) -- node[pc label, above=1pt, pos=0.35]{影响} (sel.north west);
  \end{tikzpicture}
  \caption{可视化设计核心概念图。连线标注概念间的关系类型：每条边都必须
  可读出"主语—关系—宾语"。}
  \label{fig:concept}
\end{figure}

图~\ref{fig:concept}~与思维导图的差别在于：概念图的连线承载语义（"决定""服务于"
"遵循"），表达的是概念间的\textbf{关系}而非\textbf{从属}——这是二者不可混用的根本原因。

\end{document}
